\chapter{Objectifs du projet}

\section{Objectif général}

L'objectif général du projet est de concevoir, développer et déployer une plateforme web souveraine d'automatisation à double finalité : d'une part, le \textbf{filtrage dynamique de données Excel avec création et synchronisation automatique en temps réel des feuilles enfants}, et d'autre part, la \textbf{génération et mise en forme de documents Word normalisés}, garantissant un niveau élevé de sécurité, de performance et d'ergonomie pour l'ensemble des agents de l'ANTIC.

\section{Objectifs spécifiques}

De manière opérationnelle, le projet poursuit les objectifs spécifiques suivants :

\begin{enumerate}
    \item \textbf{Élimination des tâches manuelles :} Supprimer les opérations fastidieuses de "Copier/Coller" et de reformatage manuel entre les tableurs et les documents Word ;
    \item \textbf{Filtrage dynamique et création de feuilles enfants :} Offrir une interface conviviale permettant de sélectionner une feuille parente, de définir la hauteur des en-têtes multi-lignes et d'appliquer un filtrage dynamique pour générer une nouvelle feuille enfant (\textit{Child Sheet}) au sein du classeur tableur (.xlsx vers .xlsx) ;
    \item \textbf{Mises à jour et synchronisation automatiques :} Répercuter automatiquement et en temps réel toute modification ou réactualisation des données de la feuille parente source sur les feuilles enfants dérivées, éliminant définitivement les erreurs de désynchronisation ;
    \item \textbf{Automatisation de la conversion Word :} Réduire le temps de génération d'un rapport Word complet (.docx) à moins de 8 secondes, quel que soit le nombre d'onglets du classeur source ;
    \item \textbf{Fidélité de la mise en page :} Transposer automatiquement chaque feuille de calcul en un chapitre distinct du document Word, en mode Paysage (\textit{Landscape}) avec ajustement dynamique à la page (\textit{AutoFit Window}) et préservation des bordures et en-têtes complexes ;
    \item \textbf{Souveraineté et sécurité des données :} Assurer un traitement 100\% local au sein de l'infrastructure de l'ANTIC, éliminant les fuites vers des convertisseurs publics en ligne et neutralisant les macros virales ;
    \item \textbf{Traçabilité et auditabilité :} Journaliser dans une base de données PostgreSQL immuable l'identité de l'opérateur, l'horodatage précis de l'action et le condensat SHA-256 du fichier source ;
    \item \textbf{Stockage contrôlé et purge :} Mettre en place un mécanisme de chiffrement au repos (AES-256) et de purge automatique après la période de rétention (7 jours) ou sur demande de l'agent ;
    \item \textbf{Expérience utilisateur moderne :} Offrir une interface Web réactive avec glisser-déposer (\textit{Drag \& Drop}), barres de progression dynamiques et sélection visuelle des sous-ensembles de données.
\end{enumerate}

\section{Résultats attendus}

À l'issue du projet, les résultats concrets suivants seront livrés et opérationnels :
\begin{itemize}
    \item Une application Web accessible sur le réseau intranet de l'ANTIC ;
    \item Un moteur de traitement backend léger Python FastAPI robuste et évolutif ;
    \item Une base de données PostgreSQL configurée pour la journalisation d'audit ;
    \item Une suite de conteneurs Docker facilitant le déploiement et la maintenance de la solution ;
    \item Un document complet d'homologation et de recette fonctionnelle.
\end{itemize}
